\cotefile{127}
\scanpage{2}
% Upper right. The vertical rule begins on the second line of the heading.
\noindent\begin{tabular}{@{}l@{}}
\quad\uline{Application de la}\\
\begin{tabular}{|l}
\uline{descente non plate} :\\
platitude d'une adhérence\\
schématique, et \unsure{lemme} de\\
Hironaka
\end{tabular}
\end{tabular}
% archive sheet number: 1
\scanpage{3}\mspage{1}
% archive sheet number: 2 (upper left)
% l. 1
\uline{Platitude d'une adhérence schématique}
\margin{$U\cap T\;\unsure{\frac{\phantom{T}}{|X|}}$}

% ll. 2–3. The i is beneath the inclusion sign.
1.1.
\diagram{\begin{tikzcd}[ampersand replacement=\&,column sep=small]
U\mathrel{\underset{i}{\subset}} \&[-1.5em] X\arrow[d,"f"] \& F\\
\& Y \&
\end{tikzcd}}
données $f\colon\inserted{X\to Y}$, $U$ ouvert \inserted{rétrocompact} de $X$, $F$ quasi-cohérent sur $X$
$G_U$ quotient de $F|U$, quasi-cohérent et \uline{plat} sur $Y$

% ll. 4–9
\uline{Pb} Trouver un quotient \inserted{\unsure{qu. coh.}} $G$ de $F$, qui \uline{prolonge}
$G_U$, \deleted{\unsure{et}} qui soit plat sur $Y$, et \deleted{tel} que
\unsure{i.e.} $G\xrightarrow{\inserted{\text{rel. à }i}}G_U$ \inserted{$[\,G\to i_*(G_U)$} soit \uline{univ. séparant} rel $/S$.
[Il \inserted{faut et} suffit pour
ceci, \deleted{\unsure{lorsque} $G$} si $f$ est loc. de présentation finie, et si
\deleted{$Z=\operatorname{Supp}G$, que} $G$ est de prés. finie sur $S$, que
pour tout $y\in Y$, $U_y$ contienne $\operatorname{Ass}G_y$].

% margin note written sideways
\margin{\deleted{On peut donc supposer $X$ plat sur $S$ (lisse sur $S$ si à l'origine il était \inserted{loc.} de type fini) et $F$ libre sur $X$ \deleted{\unsure{coh.}} (et de type fini s'il l'était initialement), quitte à \unsure{coiffer|coffer}}}

% ll. 10–15. The lower arrows belong to the image formula.
Si le Pb a une solution, elle est unique,
savoir
\diagram{\begin{tikzcd}[ampersand replacement=\&,column sep=small,row sep=small]
G\simeq\operatorname{Im}\bigl( F \arrow[r] \arrow[dr] \& i_*(G_U)\bigr).\\
\& i_*(F_U)\arrow[u]
\end{tikzcd}}
La question est donc de savoir si le $G$ défini par cette formule
est \deleted{\unsure{univ.}}\ \deleted{\unsure{univ.}} plat rel $/Y$ et $G\to G_U$ est
univ. \deleted{\unsure{injec.}} séparant rel $/Y$. \inserted{\uline{Remarque} Question \uline{locale} sur $X$ et $Y$, qu'on peut supposer affines !}
L'unicité implique
aussi

% ll. 16–18
\uline{Lemme} 1.2. Question locale \inserted{\uline{sur $Y$}} pour fpqc. De \unsure{m} la
\margin{SGA 1 \uline{VIII}}
question analogue, où on veut que $G$ soit \deleted{loc} de
prés. finie sur $X$. // Plus gén., descente possible \uline{sur $X$} pour morphisme fid. plat \unsure{qu. cpt|qu. cpct} $X'\to X$.

% ll. 19–20
La descente non plate \inserted{(du quotient \uline{plat} d'un faisceau déjà donné),} permet ensuite de prouver :

% ll. 21–23
\uline{Lemme} 1.3. Supposons $Y$ artinien \inserted{$=\operatorname{Spec}A$}, \deleted{\unsure{q}} et qu'il existe un
épimorphisme \inserted{$A\hookrightarrow\Gamma(Y,\mathcal O_Y)$} de schémas $Y'\to Y$, tel que après ce
changement de base, le Pb ait une solution. Supposons
% slanted
\margin{Il faut supposer soit $Y'\to Y$ fini ou $f\colon X\to Y$ loc. de \unsure{p.f.|prés. f.}, et on cherche des solutions de prés. finie}
\scanpage{4}\mspage{2}
% ll. 1–4
$Y'\to Y$ fini, \uline{ou} $X$ loc. de prés. finie, et la solution
$G'_{\deleted{\unsure{U}}}$ de présentation finie. Alors le Pb a une solution,
\deleted{sur $Y$}, et dans le $2^{\text{ème}}$ cas, la solution $G$ est aussi
de présentation finie.

\margin{$T=\overline{\{x'\}}$\par $(T_y\cap U)\cup(T_y\cap Z)=T_y$}
\margin{en les pts de $X-U$}

% ll. 5–9
\noindent\fbox{\inserted{Théorème} \deleted{Lemme} 1.4.}\quad\inserted{\uline{Néc. Suff.}}
Supposons que \deleted{\unsure{pour tt}} \inserted{$X$,} $Y$ \deleted{soit} soient loc. noeth,
les anneaux locaux de $X$ à \uline{fibres formelles} géom. \deleted{\unsure{unib.}}\ \deleted{\unsure{irréductibles}} \inserted{\unsure{loc. intègres}}.
\margin{$F$, $G$ \unsure{cohérent|cohérents}}
Supposons que \textcircled{a} pour tout couple de points $y\ne y'$ de $Y$, avec $y\in\overline{\{y'\}}$,
\deleted{\unsure{et}} tout élément \deleted{$x\in U_y$, tel q} \inserted{$x\in X_y-U_y$, et \deleted{\unsure{toute}} générisation} $x'\in\operatorname{Ass}(G_U\mid\deleted{\unsure{y}}Y_{y'})$ \inserted{de $x$}, il
existe une spécialisation $x_1$ de $x'$ \deleted{qui} qui est dans $U_y$,
et qui générise $x$, -- et que la \unsure{m\textsuperscript{me}|m\textsuperscript{e}} condition soit vérifiée sur $\forall$ $X'$ étale sur $X$.

% Left margin alongside condition (a).
\diagram{\begin{tikzcd}[ampersand replacement=\&,column sep=small,row sep=small]
U_y\ni x_1\arrow[d]\&\\
\unsure{X_y-U_y}\ni x\&x'\arrow[l]\arrow[ul]
\end{tikzcd}}
\diagram{\begin{tikzcd}[ampersand replacement=\&,column sep=small]
\underset{y}{\cdot}\&\underset{y'}{\bullet}\arrow[l]
\end{tikzcd}}

% ll. 10–16
Supposons de plus que \textcircled{b} pour $\forall$ $y\in Y$, il existe une
famille de morphismes locaux $Y_i\to Y'=\operatorname{Spec}(\widehat{\mathcal O}_{Y,y})$,
les $Y_i$ noeth locaux,
tels que après changement de base le Pb ait une
solution \deleted{de type fini} \inserted{cohérente}, \deleted{et que les \unsure{``intersections''} \unsure{finies} \unsure{des}} \inserted{$\bigcap_i\operatorname{Ker}(\widehat{\mathcal O}_{Y,y}\to\mathcal O_{Y_i,y_i})=\deleted{\unsure{\mathfrak a}}0$.}
On suppose de plus, \textcircled{c} ou bien $f$
loc. \deleted{de} de t. fini, ou \inserted{ou} bien $Y_i\to Y'$ finis. Alors le
Pb a une solution, qui est un Module cohérent
$G$ (plat sur $Y$).

% margin note written sideways
\margin{i.e. $\forall\,x'\in\operatorname{Ass}G_{U'}$, posant $T=\overline{\{x'\}}$, et pour tout $y\in Y$, $T_y\cap U$ est \uline{dense} dans \uline{$T_y$}}

% ll. 16–18
\margin{\uline{NB} Si $X$ loc. de type fini $/S$, alors la cond. a) et b) est \uline{nécessaire} pour l'existence d'une solution. a) résulte de IV 12.1.1.5. Dans ce cas, les anneaux locaux de $\uline{X}$ \unsure{d'ailleurs}, \deleted{\unsure{la}} condition quitte que sur ceux de $Y$}
\unsure{Choisir|Utiliser} pour b), et pour \unsure{v.|c)} \unsure{probant|utiliser} que les conditions sur \unsure{Lemme} 1.3. \unsure{on voit qu'en fait}
pour $Y_i$ les $\operatorname{Spec}(\mathcal O_{\deleted{\unsure{y}}Y,y}/\mathfrak m_y^{i+1})$, $i\in\mathbb N$.

% slanted
\margin{\uline{Remarque} a) est vérifiée si $\exists$ solution après tout chgt. de base $S'\to S$. \unsure{Évident}}
% slanted
\margin{du complémentaire du plus grand ouvert au dessus duquel $G$ a les propriétés qu'on veut, on peut supposer $X\to Y$ local, $U=X-\{x\}$ ($x$ point fermé de $X$).}

% ll. 19–20
\deleted{On se localise en un} En se localisant en un
point maximal de \deleted{l'image dans $Y$ de} l'ens des points $x$ \inserted{$\in X-U$} de
\deleted{\unsure{où}} qui admettent un gén. dans $X_y-U_y$ ($y=f(x)$) qui est ds $\operatorname{Ass}G_y$,
\inserted{$G=\operatorname{Im}(U^{\unsure{-}}\to i_*(G_U))$}
% ll. 21–25. Several cancellation strokes cross and surround this draft.
\deleted{qui est \unsure{pas plat}, on peut supposer que $X\to Y$ est local, et que sur $X-\{y\}$, $G$ est plat $/Y$.}
\deleted{et que $G$ est plat $/Y$ en les pts de $X-x$ \unsure{tel de locali} i.e. sauf le pt fermé de $X$. De plus, par descente plat, \unsure{point de} un cas un pt \unsure{centralisé} $x$ de $X_y$ plus \unsure{fin}, dans la \unsure{fa} de non platitude, on peut, supposer que, \unsure{par} $x\in X_y$, tel $X$ est local, et que}
On peut supposer $U=X-\{x\}$.
\deleted{[NB. $\operatorname{Ass}G_y=\operatorname{Ass}G_U\unsure{\}\,=|]\,=}\operatorname{Ass}G$]}
% archive sheet number: 3
\scanpage{5}\mspage{3}
% ll. 1–3
\inserted{Par descente plate,} on peut supposer $X$ local \deleted{complet} \inserted{hensélien}.
\inserted{\unsure{car|on} n'aura pas \unsure{utilisé|seul.}.} \deleted{(gr. \unsure{au} \unsure{m}) \unsure{l'existence} des solutions $G$}
\deleted{Que l'hypothèse sur les Ass est conservée par passage}

% Upper left, composite marginal sketch. The short undirected connectors have no head.
\diagram{\begin{tikzcd}[ampersand replacement=\&,row sep=tiny,column sep=small]
G_{U_n}\& {} \& {} \& {} \& {}\\
U_n\arrow[d,phantom,"\cap" description]\&x_1\arrow[d]\&x'\arrow[l]\arrow[dl]\&\widehat{x}_1\arrow[dr]\arrow[d]\&\widehat{x}'\in\operatorname{Ass}\widehat U\;\deleted{\unsure{\text{générisant}}}\arrow[d]\\
{}\&x\&{}\&\deleted{\unsure{x}}\&\widehat x\\
X_n\arrow[r,dash]\arrow[d,dash,"f_n"]\&X\arrow[d,dash]\&\widehat X=X'\supset U'\arrow[l]\&{}\&{}\\
Y_n\&Y\arrow[l]\&{}\&{}\&{}\\
{}\&y\&y'\arrow[l,dash]\&{}\&{}
\end{tikzcd}}

% ll. 3–8
détruit les hypothèses \inserted{\textcircled{a} \textcircled{b}} faites au départ, \deleted{\unsure{mais}} \inserted{\unsure{ni|en}}
\deleted{(quant $X^h\colon\operatorname{Spec}B^h$) on aura gagné (puisque $B^h$}
l'hypothèse que $B=\mathcal O_{X,x}$ est à fibres formelles
géom\textsuperscript{t}. normales (ce qui est en effet une
condition stable par hensélisation). Par contre,
on aura gagné que le morphisme $X'=\operatorname{Spec}\widehat B\to X=\operatorname{Spec}B$ est
à \uline{fibres formelles \inserted{(géom)} intègres}, donc si $Z_\alpha$ sont les
cycles associés à $G$ sur $X$, \deleted{et \unsure{à} $X'$} alors \deleted{\unsure{el}} ceux
associés à $\deleted{\unsure{F}}G'$ sur $X'$ sont les $Z'_\alpha$ (qui sont,
encore intègres !). Donc l'hyp. \textcircled{a}, \inserted{savoir $Z'_\alpha\cap U_y\ne\varnothing$ pr $\forall\,\alpha$,} reste alors vérifiée
en remplaçant $(X,G,U)$ par $(X',G',U')$. Cela nous
permet de remplacer $X$ par $X'$, (l'hyp. \textcircled{b} reste
trivialement vérifiée).

% ll. 15–end
Utilisant le \unsure{syst.} des solutions
du Pb obtenu pour les $Y_n$, \deleted{et} passant à la
limite, on trouve donc un module quotient
$G\to H$ de $G$, qui est plat sur $Y$, \deleted{\unsure{et}} et qui est
tel que l'\unsure{hom|homo|homom.} $G\to H$ soit \deleted{\unsure{un}} un isom sur les
pts de $U_y$ ; comme pour $\forall\,Z_\alpha$, on a $Z_\alpha\cap U_y\ne\varnothing$,
on conclut que $G\to H$ est injectif sur $U$, \deleted{donc}
injectif puisque $G\to i_*(G_U)$ est injectif. Donc
$G=H$, i.e. $G$ est \uline{plat} sur $Y$, en \unsure{$=$} temps on
voit que $G_y=H_y$ est tel que $x\notin\operatorname{Ass}G_y$, donc
on a gagné.

% Lower left margin, braced together.
\diagram{\begin{tikzcd}[ampersand replacement=\&,row sep=small,column sep=small]
F_n\arrow[r,"u_n"]\arrow[dr,"\text{surj}"']\&i_{n*}(G_{nU})\\
{}\&G_n\arrow[u]\arrow[r,"\unsure{v_n}"']\&H_n
\end{tikzcd}}
\margin{$H_n=\operatorname{Im}u_n=\operatorname{Im}v_n$ est un quotient de $G_n$}
% archive sheet number: 4
\scanpage{6}\mspage{4}
% ll. 1–10. A light pencil X runs through the first corollary.
\deleted{\uline{Cor.} 1.5. Supposons que $f$ soit loc. de type fini, \inserted{$F,G$ cohérents} et que
\textcircled{a'} les $(G_U)_y$ soient \inserted{$(S_1)$ et} équidimensionnels de dim. donnée $d$,
tandis que (définissant encore $G$ comme $\operatorname{Im}(F\to i_*(G_U))$]
\deleted{que $U$ soit \unsure{schi de|schém. dense}} pour $\forall\,y$, $U_y$ soit dense dans
$X_y$. // \deleted{Alors} Supposons de plus que pour $x\in\deleted{\unsure{X_U}\unsure{G}}\inserted{X-U}$,
l'anneau local $\mathcal O_{X,x}$ soit à fibres formelles géom
normales \deleted{Alors la condition} (il suffit qu'il en soit de
\unsure{m.} de $\mathcal O_{Y,f(x)}$ !). Alors la condition \textcircled{b} du th.
est néc. et suffisante pour l'existence d'une solution
du Pb.}
% slanted
\margin{\uline{$\operatorname{Supp}G=X$ ?}}
\margin{\deleted{\illegible{1}}}

% ll. 11–16
\uline{Critère valuatif} :

\uline{Cor.} 1.6. Notations \deleted{et} hyp. préliminaires de 1.4. \inserted{y compris l'hyp. des fibres formelles}
On suppose de
plus que pour $\forall\,x\in X-U$, posant $y=f(x)$, $\mathcal O_{Y,y}$ \uline{soit réduit}, \deleted{et}
\deleted{\unsure{$Y$ soit soit}}, et que pour $\forall$ changement de base $Y'\to Y$,
avec $Y'$ un \deleted{\unsure{dvr trait}} \inserted{trait} \deleted{\unsure{fini sur $\mathcal O_{Y,y}$}}
\deleted{\unsure{spectre d'un anneau}} ($=$ spectre d'un ann. de val
\unsure{discrèt|discrète}) \inserted{fini sur $\operatorname{Spec}\widehat{\mathcal O}_{Y,y}$,} on ait une solution du Pb. Alors on
a une solution du Pb.
% slanted
\margin{1.5 ? \unsure{sinon}, manque apparemment hypothèse du type de 1.4 a)}
% slanted
\margin{? Fait double emploi avec 1.7, sauf qu'on \unsure{en|ne|a} fait d'hyp. ``de type fini'' \unsure{--} ?}

% ll. 17–23
\uline{Cor.} 1.7. $f\colon X\to Y$ \deleted{loc. de type fini}, \deleted{\illegible{1}} \inserted{$Y$ loc. noeth,} $F,G_U$ \unsure{cohérent|cohérents} \deleted{\unsure{plats}}
\deleted{\unsure{munis d'intentions}}. On suppose :

a) Pour $\forall\,x\in X-U$, posant $y=f(x)$, $\mathcal O_{Y,y}$ est \deleted{\unsure{intègre}}
réduit, et à fibres \deleted{\unsure{géom.}} formelles géom. \deleted{\unsure{irréductibles}}, \deleted{\unsure{unibranches}} \inserted{normales}
(p. ex. an. de type fini sur $\operatorname{Spec}\mathbb Z$).

b) Pour $y$ comme dessus, et tout \deleted{de type \unsure{fini}} $Y'$ fini sur
$\operatorname{Spec}\mathcal O_{Y,y}$, il y a une solution du Pb après chgt de
base $Y'\to Y$.
\margin{\uline{NB} Dans cette condition, l'assertion de platitude devient triviale, seule la question \deleted{\unsure{de|des}} $\operatorname{Ass}G'_{y'}\subset U'_{y'}$ est sérieuse}
\margin{\uline{\unsure{Expliciter|Explicitez}}}

\deleted{c) La condition a) de 1.4 est satisfaite, ou la condition a') de 1.5.}
% archive sheet number: 5
\scanpage{7}\mspage{5}
\uline{Alors} le Pb posé admet une solution.

Nous \unsure{allons} maintenant utiliser le \uline{lemme de Hironaka}.
% archive sheet number: 6
\scanpage{8}\mspage{6}
\uline{Th. 1.8.}

$f\colon X\longrightarrow Y$ loc. de prés. finie

\margin{\uline{NB} Si on suppose $Y$ loc. noeth, alors l'énoncé équivaut à dire que, si $X$ est réduit, alors au voisinage de $x$ $f$ est normal.}
\margin{Si $X$ et $Y$ sont loc. noeth, on doit pouvoir supprimer l'hyp. que $f$ est \deleted{\unsure{loc.} \unsure{de t.}} de prés finie (moyennant des hypothèses d'excellence éventuellement --}

\deleted{\unsure{\textcircled{1}}} \unsure{Soient|Soit} $x\in X$, $y=f(x)$

a) $Y$ réduit en $y$

b) $f$ univ. ouvert en les gén. maximales de $x$ dans $X_y$

c) Pour tte gén. \deleted{\unsure{maximale}} $y'$ de $y$, et toute \inserted{les} générisation $x'$ de
$x$ appartenant à $X_{y'}$, les $\dim_{x'}X_{y'}$ sont égaux à \deleted{\unsure{$\dim X_y$}} $\dim_x\deleted{\unsure{X}}X_y$

d) $X_y$ \deleted{\unsure{est|réduit} \unsure{$x$} séparable en} réduit en les gén. maximales de $x$ dans $X_y$

e) $(X_y)_{\text{réd}}$ est géom\textsuperscript{t}. normal au voisinage de $x$

% slanted
\margin{peut être remplacé par c') $f$ univ. ouvert dans un voisinage de $x$, et c) \unsure{valable si} $y'$ est maximal}

[d) et e) équivalent à : $\exists$ voisinage ouvert $U$ de $x$ \inserted{dans $X_y$}, tel que
$\overline U$ \inserted{$=U\otimes_{k(y)}\overline{k(y)}$} soit gén. réduit et $\overline U_{\text{réd}}$ normal].

% Conclusion and abandoned formulations, from the centre down.
\uline{Alors} \deleted{\unsure{1°|1}} Il existe un voisinage ouvert $U$ de $x$, tel que
le morphisme \inserted{$U_{\text{réd}}\to Y$ soit \deleted{normal} loc. de présentation finie et}
\deleted{$U\to Y$ soit univ. ouvert, \unsure{et} \unsure{intègre} \unsure{fini} et à fibres géométriques normales, donc \inserted{localement} équidimensionnel}
\deleted{à fibres géométriques \inserted{\unsure{gén.}} réduites telles que $(X_y)_{\text{réd}}$ normal.}

\deleted{2) Si $Y$ est réduit, la condition envisagée est vérifiée,}
\deleted{et si $Y$ est réduit, alors \quad Soit $Y_0=Y_{\text{réd}}$, $X_0=X\times_Y Y_0$.}

\deleted{\uline{Cor.} \textcircled{?} Conditions équivalentes : (i) Les conditions envisagées dans \textcircled{1} sont satisfaites en $\forall\,x\in X$, \unsure{et que $f$ est univ. ouvert} fibre \unsure{géométriquement} \unsure{unib|univ.} \inserted{i.e. (i bis) $f$ est univ. ouvert, \inserted{loc} équidimensionnel, ses fibres sont gén. séparables et ses fibres réduites géom. normales.}}
\deleted{(ii) $\exists$ un sous--}
\deleted{(ii) Il existe un ouvert $V$ de $X_0$, tel que pour $\forall\,y\in Y$, $V_y$ soit dense dans $X_y$, et $V\longrightarrow Y_0$ un morphisme \deleted{normal} séparable. De plus \deleted{$X_0\to Y_0$ est à fibres loc. équidimensionnelles, et} pour $\forall\,y\in Y_0$, $(X_{0y})_{\text{réd}}$ est géom. normal.}
\deleted{(iii) \deleted{Il existe un sous $X_{0\,\text{réd}}\to Y_0$} $X_{\text{réd}}\to Y_{\text{réd}}$ est normal et \inserted{loc.} de présentation finie}
\deleted{(iiii bis) Il existe un sous-schéma \deleted{\unsure{X}} $Z_0$ de $X_0$, \deleted{plat} normal \inserted{et loc. de prés. finie} sur $Y_0$, \deleted{et} ayant \unsure{m\textsuperscript{me}} espace sous-jacent que $X_0$, et coïncidant avec $X_0$ en les pts génériques des fibres \ldots}

normal. De plus, il existe un ouvert \inserted{relativement} schématiquement dense $V$ de
$U$ \deleted{\unsure{du|dans}} $S$, tel que $V\to X$ soit une \unsure{imm.|univ.} ouverte
% archive sheet number: 7
\scanpage{9}\mspage{- IV-\handwritten{1131} -}
% Single long, near-vertical pencil stroke through the setting copy.
\margin{\handwritten{\unsure{?}}}
existe , en vertu de (21.9.12) un $\underline{\mathrm O}_X$-Module inversible $\underline L$ dont \deleted{\illegible{1}}
l'image réciproque dans $X'_0$ est isomorphe à $\underline L_0$ , donc est ample . Puisque tout voisinage de $s_0$ dans $S$ est nécessairement $S$ tout entier ,on déduit alors de (9.6.4) que $\underline L$ est un $\underline{\mathrm O}_X$-Module ample , d'où la conclusion \deleted{\illegible{1}} (\uline{II},5.3.1 et \uline{II},5.5.3)\deleted{\illegible{1}} .

\uline{21.10 . Factorialité des anneaux locaux réguliers} .

\uline{Théorème} (21.10.1) \inserted{(Auslander-Buchsbaum)} .- \uline{Un anneau local noethérien régulier est factorie\handwritten{l}}.
% slanted
\margin{\handwritten{Créditer Kaplansky de la \unsure{démonstration|démonstr.} donnée ici.}}

Soit $A$ un anneau local noethérien régulier de dimension $n$ ; nous allon\handwritten{s} raisonner par récurrence sur $n$ . Pour $n=\unsure{0}$,$A$ est un corps et pour $n=1$ ,
un anneau de valuation discrète , donc principal (et \uline{a fortiori} factoriel) . Supposons donc $n\handwritten{\geqslant}2$ et le théorème démontré pour des anneaux réguliers d\handwritten{e} dimension $\handwritten{<}n$ . Posons $X=\operatorname{Spec}(A)$ , désignons par $a$ le point fermé de $X$ et posons $U=X-\{a\}$ . En tout point $y\in U$ , on a \deleted{\illegible{1}},
$\dim(\underline{\mathrm O}_{X,y})\handwritten{\leqslant}n-1$ , et puisque $A$ est régulier , les anneaux $\underline{\mathrm O}_{X,y}$ le sont aussi (\uline{0},17.3.2), donc l'hypothèse de récurrence entraîne qu'ils sont factoriels . De plus on a $\operatorname{prof}(A)=\dim(A)\handwritten{\geqslant}2$ puisque $A$ est régulier , donc de Cohen-Macaulay (\uline{0},17.1.3) . Utilisant (21.6.14) , on est ramen\handwritten{é} à prouver que $\operatorname{Pic}(U)=0$ . Considérons donc un $\underline{\mathrm O}_U$-Module inversible $\underline L$ , et prouvons qu'il est isomorphe à $\underline{\mathrm O}_U$ . \deleted{Puisque $\dim(A)\handwritten{\geqslant}2$ et que $A$ est régulier ,} \handwritten{I}l résulte de \deleted{(5.11.6)} \inserted{\handwritten{I 9.4.5.}} qu'il existe un $\underline{\mathrm O}_X$-Module \uline{cohérent} $\underline F$ tel que $\underline F\handwritten{\mid}U=\underline L$ . Comme $A$ est régulier , donc de dimension cohomologique finie (\uline{0},17.3.1) , il existe une réso\handwritten{l}ution gauche \uline{finie} de $F$ :
\[
0\handwritten{\longleftarrow}\underline F\handwritten{\longleftarrow}\underline{\mathrm O}_X^{n_1}\handwritten{\longleftarrow}\underline{\mathrm O}_X^{n_2}\handwritten{\longleftarrow}\mathinner{\ldotp\ldotp\ldotp\ldotp}\handwritten{\longleftarrow}\underline{\mathrm O}_X^{n_h}\handwritten{\longleftarrow}0
\]
(\uline{0},17.2.8 et \uline{0},17.2.2,(iii)) . Par restriction à $U$ , on a donc une résolution finie
\begingroup\makeatletter\tagsleft@true\makeatother
\[
0\handwritten{\longleftarrow}\underline L\handwritten{\longleftarrow}\underline{\mathrm O}_U^{n_1}\handwritten{\longleftarrow}\underline{\mathrm O}_U^{n_2}\handwritten{\longleftarrow}\mathinner{\ldotp\ldotp\ldotp\ldotp}\handwritten{\longleftarrow}\underline{\mathrm O}_U^{n_h}\handwritten{\longleftarrow}0
\tag*{\handwritten{(21.10.1.1)}}
\]
\endgroup

Le théorème résultera alors des considérations générales suivantes \handwritten{\unsure{;|:}}

Sur un espace \deleted{\illegible{1}} annelé $X$ , soit $\underline E$ un $\underline{\mathrm O}_X$ -Module localement libre
% archive sheet number: 8
\scanpage{10}\mspage{\deleted{\unsure{6|8}}7}
% ll. 1–10
\uline{Dém} \deleted{On est ramené par \unsure{principe|procédé} standard} \inserted{La conclusion \inserted{sur $U$} signifie aussi qu'}il existe
un sous-préschéma $U_0$ \deleted{\unsure{$Z$}} de $U$, de prés. finie sur $Y_0$,
tel que $U_0\to U$ soit \uline{surjectif}, \deleted{et tel que} tel que
$U_0\to\deleted{\unsure{X}}Y$ \deleted{soit tel}, soit \deleted{plat} normal \deleted{à \unsure{géom.} \unsure{et} donc loc. de prés finie}
(ce qui implique donc que $U_0$ est réduit, donc $U_0=U_{\text{red}}$), et un ouvert
$V$ de $U_0$, schématiquement dense rel $/Y$, tel que
$V\longrightarrow X$ soit une imm ouverte. Sous cette
forme, on voit qu'on est ramené au cas
où $Y$ est le spectre d'un anneau local ess.
de type fini sur $\mathbb Z$.
\deleted{A prouver qu'il existe une \illegible{1}}

% ll. 11–16
Utilisant la technique \inserted{de constructibilité et} de passage à un \unsure{sch. ét.|échelon|limite},
et le lemme de Hironaka, on voit que l'hypothèse sur la fibre $X_y$ en $x$ est vraie encore aux
pts voisins \inserted{avec la \unsure{m\textsuperscript{me}} dim locale $d$.} Utilisant \inserted{EGA \uline{IV}} 16.1. \deleted{\unsure{et sa}} et la \deleted{\unsure{3\textsuperscript{ème}}}
condition \textcircled{c}, on trouve que $f$ est univ. ouverte \inserted{et équidim} dans
un voisinage de $x$. \deleted{D'autre part, les \unsure{conditions}}
\inserted{On peut donc supposer $f$ univ. ouvert, équidim, à fibres géom. gén. réduites,}
et à \deleted{fibres géom. \unsure{chacunes}} telles que $(\overline X_y)_{\text{réd}}$ soit normal. \deleted{\unsure{Ainsi}} On
\deleted{\unsure{15.3}, on sait que} peut supposer $X$ réduit : prouvons que
$X\to Y$ est normal.
% slanted
\margin{utiliser \textcircled{b} \textcircled{c} \textcircled{d} \textcircled{e}}
% slanted
\margin{\uline{Lemme à \unsure{dégager|dégg.er}.}}
% slanted
\margin{\deleted{On peut \unsure{procéder} par récurrence sur $\dim Y$, et supposer que l'on \unsure{gagne} en dim plus petite}}

% ll. 17–21
Soit $V$ l'ouvert des pts $x$ de $X$ tels que $X_{y=f(x)}$ soit
\uline{réduit} (donc séparable) en $x$ sur $k(y)$. En vertu de \textcircled{a} et
\inserted{EGA \uline{IV}} 15.2.3. c'est aussi le plus grand ouvert sur lequel
$f$ induit un morphisme \deleted{réduit} séparable, et en
vertu de l'hypothèse faite (déduite de \textcircled{d,e}) $V$ est
dense dans $X$ fibre par fibre. \deleted{Si donc $X$ est normal aux \unsure{pts fermés} \unsure{$X_{\text{réd}}$ a les} parties \unsure{réduites},}

% ll. 22–24. Intersecting rulings and text running in opposite directions.
\deleted{Si $\overline X$ est l'adhérence schématique \deleted{\unsure{il sera}} de $V$. Conclusion (que $X$ \inserted{et $Y$} est réduit (qui \unsure{va} se perdre par ext. de la base), \deleted{mais}}
% margin note written sideways (inverted)
\margin{\deleted{\uline{Téléphoner à Coates}}}
% margin note written sideways (inverted)
\margin{\deleted{\unsure{Jean Artin} \unsure{pour 4 heures|pour 1 heure}, voir \unsure{Feitings|fibrés}, \unsure{batterie} \unsure{un} \unsure{cas} \unsure{lost} \unsure{letter}}}

\margin{1.7 ci-dessus}
Le critère valuatif \unsure{nous} ramène, lui, au cas où $Y$ est le sp. d'un anneau de val. discrète,
(par récurr\unsure{ence} !) c'est le cas de Hironaka, OK.
% archive sheet number: 9
\scanpage{11}
\diagram{\begin{tikzcd}[ampersand replacement=\&,column sep=small]
x\&X\arrow[d,"f\ \text{de t.f.}"]\&{}\\
y\&Y\&\text{\uline{noeth}}
\end{tikzcd}}
\[
\left\{\begin{aligned}
&\mathcal O_{Y,y}\ \text{réduit}\\
&\mathcal O_{X,x}\ \text{réduit}\\
&\exists\,d\ \text{tel que, au voisinage de }x,\ \text{ttes les comp. des fibres de } f\\[-2pt]
&\hspace{12em}\text{sont de dim }d\\
& f\ \text{universellement ouvert}\qquad(\text{EGA }\underline{\mathrm{IV}}\ 14,15)\\
&\left\{\begin{aligned}
&X_y\quad\text{gén. maximales }\inserted{x_i}/\text{de }x\text{ ds }X_y\ ;\ X_y\ \text{réduit en les }x_i\\
&(X_y)_{\text{red}}\ \text{est géom normal }/k(y)
\end{aligned}\right.
\end{aligned}\right.
\]
\[
\begin{aligned}
\Longleftrightarrow\quad &\overline{X_y}\ (\overset{\text{dfn}}{=}X_y\otimes_{k(y)}\overline{k(y)})\ \text{est, au voisinage}\\
&\text{de }\bar x,\ \text{gén. réduit, et }(\overline{X_y})_{\text{red}}\ \text{est normal}\\
\Downarrow\qquad&\\
\exists\quad&\left|\begin{aligned}
&\text{voisinage ouvert }U\text{ de }X\text{ tel que }U\to Y\\
&\text{soit ``normal'' i.e. \uline{plat} à fibres géom\textsuperscript{t}.}\\
&\text{normales.}
\end{aligned}\right.
\end{aligned}
\]

\uline{Dém} \uline{lorsque on suppose} $\mathcal O_{Y,y}$ \uline{excellent} [il suffit : à fibres formelles géom. normales]

a) Ttes les hyp. \deleted{\unsure{sont}} \inserted{restent} vraies au voisinage \inserted{$U$} de \deleted{\unsure{$x$}}

EGA \uline{IV} 9
\diagram{\begin{tikzcd}[ampersand replacement=\&,column sep=small,row sep=small]
x_0\quad x'_0\&X'\arrow[r,dash]\arrow[d,dash]\&X\arrow[d,dash]\&x\quad x'\\
{}\&Y'\arrow[r,dash]\arrow[u,shift right]\arrow[ur]\&Y\&y\quad y'\\
{}\& y_0\quad y'_0\&{}\&{}
\end{tikzcd}}
\diagram{\begin{tikzcd}[ampersand replacement=\&,column sep=small,row sep=small]
{}\&{}\&X\arrow[dl,dash]\arrow[dd,dash]\&x\ \unsure{,}\ x'\\
x'_0\quad x_0''\&X''\arrow[dd,dash]\&{}\&{}\\
{}\&{}\&Y\&y-y'\\
{}\&Y'\arrow[ur,dash]\arrow[uu,dashed,shift right]\arrow[uur,dashed]\&{}\&{}\\
{}\&y''_0\quad y_0\&{}\&{}
\end{tikzcd}}

\uline{vérifions} l'ens des pts de $X'$ où $X'$ est réduit est \unsure{vrai|vraie|avec} sur chaque fibre,

\unsure{Remarquer|Remarque|Remarquons} que l'ens \inserted{ouvert $V$} des pts où $X/Y$ est \deleted{plat} ``séparable'' i.e. \uline{plat} à fibres séparables est, \deleted{\unsure{au}} voisinage de $x$, dense dans chaque fibre.

$V$ est un ouvert : EGA \uline{IV} 12

$V$ contient les pts de $X_y$ où $X_y$ est réduit
\[
\text{EGA }\underline{\mathrm{IV}}\quad \unsure{15\text{r}|15.}2.3
\]
\margin{donc $V\cap X_y$ est dense dans $X_y$, et on conclut par  équidimensionnalité}
% archive sheet number: 10
\scanpage{12}
\uline{NB} \deleted{\deleted{\unsure{Soit}} \unsure{Supp.|Supposons} $S=\operatorname{Spec}k$, \deleted{\unsure{$X=$}} $k$ corps alg. \uline{clos}, et $X=S$,}
\[
\begin{aligned}
\deleted{\unsure{S_0}}&\quad\deleted{S_0=\text{schémas \inserted{de t.f.} étales sur }X=\text{\unsure{isom.} locaux sur }X,}\\
&\quad\deleted{\mathcal S=\text{schémas de t.f. sur }S}\\
&\quad\deleted{\mathcal M=\text{imm. ouvertes}}
\end{aligned}
\]

\deleted{La catégorie $\mathcal C$ est celle des schémas \deleted{\unsure{ét.}} finis, sur $k$, avec la top. de Zariski, un faisceau \inserted{de $\mathcal O$-modules} dessus revient à la donnée d'un \deleted{module} foncteur $\Lambda\leadsto F(\Lambda)$, $\Lambda$ fini sur $k$, commutant aux produits.}

\deleted{Les faisceaux qu. coh. \inserted{sur $\operatorname{Spec}\Lambda$} sont ceux \unsure{pour} \deleted{\unsure{lesquels}} de la forme $\deleted{\unsure{F}}\ \Lambda'\leadsto M\otimes_{\Lambda}\Lambda'$, avec $M$ un $\Lambda$-module, les faisceaux de prés. finie sont ceux pour lesquels $\deleted{\dim_k}M$ \deleted{\unsure{finie}} \inserted{de t.f.}. Mais $\mathcal O$ \unsure{lui-m.} n'est pas cohérent, car \deleted{\illegible{1}} soit un hom $\mathcal O_\Lambda^n\overset u{-}\mathcal O_\Lambda$ défini par \deleted{\illegible{1}} $u_1\ldots u_n\in\Lambda$, alors $\operatorname{Ker}u$ est un ``faisceau sur \deleted{Spec} $\Lambda$'' \unsure{f.|?}}
\[
\deleted{\Lambda'\longmapsto\operatorname{Ker}\bigl(u_{\Lambda'}:\deleted{\unsure{\mathcal O_\Lambda}}\Lambda'{}^n-\Lambda'\bigr)}
\]
\deleted{qui n'est pas en général déduit de $\operatorname{Ker}u$ comme $\Lambda'\leadsto(\operatorname{Ker}u)\otimes_\Lambda\Lambda'$, manque de platitude \ldots\,.}

% margin note written sideways
\margin{\diagram{\begin{tikzcd}[ampersand replacement=\&,column sep=small]
T'_y\arrow[d,dash] \& T'\arrow[d,"\unsure{\text{ét.}}"'] \& X'\arrow[d] \& x'\\
T_y \& T \& X \& x
\end{tikzcd}}
$U\cap T_y$ dense dans $T_y$\\
$U'\cap T'_y$ dense \unsure{ds|dans} $T'_y$\\
$T'_{\unsure{k|i}}=\bigcup T'_i$\\
$T_y=\bigcup T'_{iy}$}
% margin note written sideways, in the opposite direction to the main text
\margin{\diagram{\begin{tikzcd}[ampersand replacement=\&,column sep=small]
T \&[-1em] \&\\
\tikz[baseline=(a.base)]{\node[inner sep=1pt](a){$\overline{\{x'\}}$};\draw (a.center) ellipse (1.0em and 1.3em);}\cap X_y \& U_y \&\\[-.5em]
\shortparallel \& \&\\[-.5em]
T_y \& \& T_{y\unsure{a}}
\end{tikzcd}}}
% archive sheet number: 11
\scanpage{13}
b) \uline{On suppose les hyp. satisfaites en tous les pts de $X$.}

$U$ l'ouvert où $f$ est séparable, il est dense sur chaque fibre. \unsure{Á|À} prouver que $X=\overline U$ \inserted{(adh. schématique)}, est \deleted{\unsure{p}} \uuline{plat} sur $Y$.

On \unsure{aurait|avait} terminé, \unsure{or|on} il reste à prouver $X_y$ est \unsure{séparable}. Par hyp., or soit $y'$ une gén. maximale de $y$, et appliquons $Y'\longrightarrow Y$, donc $X'_{y'_0}$ est réduit, donc $X'$ (plat sur $Y'$, donc $=\overline{X'_{y'_0}}$ \inserted{(adh. sch.)}) est aussi réduit. Donc théorème s'applique, $X'_{y_0}$ \unsure{''} est \unsure{normale|normal}, donc $X_y$ aussi \uline{cqfd}

c) Pour prouver la platitude de $X/Y$, on applique le \uline{critère valuatif de platitude et \unsure{aux|une} adhérences schématiques}
% archive sheet number: 12
